\chapter{Glucose Tolerance Test (OGTT)}

\section*{What Is This Test?}
An oral glucose tolerance test measures how the body processes a known glucose load. Blood glucose is measured while fasting and again at specified times after the patient drinks a glucose solution. It is used in selected diabetes evaluations and is a key test for diagnosing gestational diabetes according to the applicable pregnancy protocol.

\section*{Alternative Names}
OGTT; Oral Glucose Tolerance Test; GTT; Glucose Challenge Follow-up; Two-Hour Glucose Tolerance Test; Three-Hour Glucose Tolerance Test

\section*{Why It Is Ordered}
Evaluation for gestational diabetes, confirmation or clarification of prediabetes or diabetes when other results are borderline, and selected assessment of abnormal glucose handling. The pregnancy and nonpregnancy protocols use different glucose loads, sampling times, and diagnostic thresholds.

\section*{How This Test Is Performed}
After an overnight fast, a baseline blood sample is collected. The patient drinks a measured glucose solution, and additional blood samples are collected at the protocol-defined times, commonly at 1 and 2 hours or over 2--3 hours. The patient should remain seated and avoid eating, smoking, or strenuous activity during the test unless instructed otherwise.

\section*{Preparation, Risks, and Limitations}
Follow the laboratory's fasting and medication instructions; do not change diabetes medicines without clinical direction. Nausea, sweating, dizziness, or temporary high or low glucose can occur. Acute illness, medication effects, inadequate fasting, vomiting, and failure to complete the timed samples can make results unreliable. Interpretation must use the correct pregnancy, age, and laboratory criteria.

\section*{References}
\begin{enumerate}
  \item MedlinePlus. Diabetes Tests and Oral Glucose Tolerance Test.
  \item American Diabetes Association. Standards of Care in Diabetes.
\end{enumerate}

\clearpage
